%% The MIT License (MIT)
%%
%% Copyright (c) 2015 Daniil Belyakov
%%
%% Permission is hereby granted, free of charge, to any person obtaining a copy
%% of this software and associated documentation files (the "Software"), to deal
%% in the Software without restriction, including without limitation the rights
%% to use, copy, modify, merge, publish, distribute, sublicense, and/or sell
%% copies of the Software, and to permit persons to whom the Software is
%% furnished to do so, subject to the following conditions:
%%
%% The above copyright notice and this permission notice shall be included in all
%% copies or substantial portions of the Software.
%%
%% THE SOFTWARE IS PROVIDED "AS IS", WITHOUT WARRANTY OF ANY KIND, EXPRESS OR
%% IMPLIED, INCLUDING BUT NOT LIMITED TO THE WARRANTIES OF MERCHANTABILITY,
%% FITNESS FOR A PARTICULAR PURPOSE AND NONINFRINGEMENT. IN NO EVENT SHALL THE
%% AUTHORS OR COPYRIGHT HOLDERS BE LIABLE FOR ANY CLAIM, DAMAGES OR OTHER
%% LIABILITY, WHETHER IN AN ACTION OF CONTRACT, TORT OR OTHERWISE, ARISING FROM,
%% OUT OF OR IN CONNECTION WITH THE SOFTWARE OR THE USE OR OTHER DEALINGS IN THE
%% SOFTWARE.

% The font could be set to Windows-specific Calibri by using the 'calibri' option
\documentclass[]{mcdowellcv}

% For mathematical symbols
\usepackage{amsmath}
\usepackage{hyperref}
\usepackage{xcolor}
\usepackage{setspace}
\setstretch{1.22} % Erhöhter Zeilenabstand für gleichmäßigen, harmonischen Textfluss

% Customize hyperlink appearance
\hypersetup{
    colorlinks=true,    % Enable colored links
    linkcolor=blue,     % Color for internal links
    urlcolor=blue,      % Color for external links (URLs)
    citecolor=blue      % Color for citations
}

% Set applicant's personal data for header (Public Version)
\name{M. Salih Sayal}
\address{Duisburg, Deutschland \linebreak \mbox{muhammet.salih.sayal@hotmail.com}}
\contacts{github.com/\href{https://github.com/salihsayal}{salihsayal} \linebreak \href{https://www.linkedin.com/in/salihsayal}{LinkedIn Profil}}

\begin{document}
    
    % Print the header
    \makeheader
    
    % =========================================================================
    % SEITE 1: PRAXIS & SYSTEMERFAHRUNG
    % =========================================================================
    
    \begin{cvsection}{Berufserfahrung}
        \begin{cvsubsection}{IT-Engineer Automation}{Everllence (ehem. MAN ES)}{06/2025 -- 09/2026} 
            \begin{itemize}
                \item \textbf{Data Engineering \& ETL:} Entwicklung einer \textbf{Docker-basierten Python-Pipeline} zur Automatisierung komplexer Abrechnungsworkflows (PostgreSQL) inkl. \textbf{Active Directory (LDAP)}-Synchronisation
                \item \textbf{Backend-Entwicklung:} Konzeption und Implementierung einer \textbf{REST-API} mit \textbf{Java Spring Boot} zur nahtlosen Anbindung an \textbf{SAP-Systeme}
                \item \textbf{System-Integration:} Entwicklung spezialisierter \textbf{Wrapper-Skripte} in Python zur Orchestrierung der Interkommunikation heterogener Software-Systeme
                \item \textbf{Cloud Automation \& Reporting:} Betrieb von Workloads auf \textbf{Microsoft Azure} (VMs, Entra ID) sowie automatisiertes Daten-Reporting via \textbf{Microsoft Graph API} (OAuth2) und Plotly
                \item \textbf{Disaster Recovery \& Scripting:} Implementierung automatisierter \textbf{Backup-/Restore-Prozesse (SFTP/Rsync)} in Linux-Serverumgebungen sowie Refactoring von \textbf{Bash-Skripten}
            \end{itemize}
        \end{cvsubsection}
        
        \vspace{0.15cm}
        
        \begin{cvsubsection}{Software Engineer}{rola Security (T-Systems)}{10/2023 -- 04/2025}
            \begin{itemize}
                \item \textbf{DevOps \& Troubleshooting:} \textbf{Fehleranalyse} beim Systemstart und Automatisierung von Konfigurationseinspielungen im Linux-Umfeld
                \item \textbf{Data Engineering:} Entwicklung \textbf{Java-basierter Analyse-Tools} zur automatisierten Auswertung komplexer System-Fehlerprotokolle
                \item \textbf{Software Quality Assurance:} Sicherstellung hoher Qualitätsstandards durch die Pflege von \textbf{Checkstyle-Regeln} und Umsetzung von \textbf{Sonar-Vorschlägen} (Clean Code)
            \end{itemize}
        \end{cvsubsection}
    \end{cvsection}

    \vspace{0.2cm}

    \begin{cvsection}{Projekte}
        \begin{cvsubsection}{Infrastruktur- \& Cloud-Projekt}{Jobentra}{07/2026 -- 08/2026}
            \begin{itemize}
                \item \textbf{Kubernetes Cluster \& GitOps:} Setup und Administration eines stateless \textbf{3-Node-Kubernetes-Clusters} auf Ubuntu Server für Microservices inklusive \textbf{GitOps-Pipelines} via \textbf{ArgoCD}, Gitea und CI/CD-Runnern
                \item \textbf{Proxmox VE \& HA-Architektur:} Konzeption einer hochverfügbaren Infrastruktur auf Proxmox VE mit \textbf{aktiver Replikation} und Ausfallsicherheit auf dedizierter Primär- und Fallback-Hardware
                \item \textbf{Stateful Services \& IAM:} Betrieb zentraler Infrastruktur-Dienste wie \textbf{Samba AD DC (SSO)}, \textbf{NFS-Storage} für Kubernetes Persistent Volumes, PostgreSQL und Nginx Reverse Proxy
                \item \textbf{Full-Stack CRM-Entwicklung:} Konzeption und KI-gestützte Implementierung einer CRM-Software mit \textbf{Java Spring}, \textbf{Python FastAPI (Vertex AI / Gemini)}, \textbf{PostgreSQL} und \textbf{Next.js}
            \end{itemize}
        \end{cvsubsection}
    \end{cvsection}

    % =========================================================================
    % SEITE 2: SKILLS, BILDUNG & SPRACHEN
    % =========================================================================
    \newpage 

    \begin{cvsection}{Technische Erfahrung}
        \begin{cvsubsection}{}{}{}
            \begin{itemize}
                \item \textbf{Sprachen \& Frameworks:} Java (Spring Boot), Python (\textbf{FastAPI}, Pandas, Plotly), SQL, JavaScript / TypeScript, \textbf{Next.js}, Bash, PowerShell
                \item \textbf{Infrastruktur, Cloud \& DevOps:} Kubernetes (K3s/K8s), ArgoCD (GitOps), Proxmox VE, Docker / Compose, \textbf{Microsoft Azure}, Linux-Administration (Debian/Ubuntu, Arch, RHEL/Fedora), ZFS, CI/CD
                \item \textbf{Security \& Networking:} Identity Management (\textbf{Active Directory, LDAP, Samba AD DC}), Nginx Reverse Proxy, OPNsense Firewall, WireGuard / Tailscale VPN, NFS Storage, SSH, TLS, DNS
                \item \textbf{Tools \& Quality:} Git (Gitea/GitHub/Bitbucket), Jira, Confluence, OpenProject, API-Integration (REST, MS Graph, SAP), JUnit 5, Mockito, SonarQube, Checkstyle
            \end{itemize}
        \end{cvsubsection}
    \end{cvsection}

    \vspace{0.25cm}

    \begin{cvsection}{Bildungsweg}
        \begin{cvsubsection}{Informatik (B.Sc.)}{HHU Düsseldorf}{vorauss. Frühjahr 2027}
            \begin{itemize}
                \item Aktueller Notendurchschnitt: \textbf{2,4} (148 ECTS)
                \item \textbf{Studienschwerpunkte:} Softwareentwicklung (Java/C), Datenbanken (SQL/NoSQL)
                \item \textbf{Besondere Leistungen:} Programmierung (Note: 1,0), Datenbanken: Weiterführende Konzepte (Note: 1,0)
                \item \textbf{Wahlvertiefung:} Verteilte Systeme, Betriebssysteme, Datenbanken \& NoSQL
            \end{itemize}
        \end{cvsubsection}
        
        \vspace{0.15cm}
        
        \begin{cvsubsection}{Freiwilliger Wehrdienst}{Bundeswehr}{02/2021 -- 06/2021}
        \end{cvsubsection}
        
        \vspace{0.15cm}
        
        \begin{cvsubsection}{Allgemeine Hochschulreife}{Elly-Heuss-Knapp-Gymnasium Duisburg}{2020}
        \end{cvsubsection}
    \end{cvsection}

    \vspace{0.25cm}

    \begin{cvsection}{Sprachkenntnisse}
        \begin{cvsubsection}{}{}{}     
            \begin{itemize}
                \item Deutsch (Muttersprache)
                \item Englisch (fließend)
                \item Türkisch (fließend)
                \item Japanisch (fortgeschrittene Grundkenntnisse)
            \end{itemize}
        \end{cvsubsection}
    \end{cvsection}
    
\end{document}